%%%%%%%%%%%%%%%%%%%%%%%%%%%%%%%%%%%%%%%%%%%%%%%%%%%%%%%%%%%%%%%%%%%%%%%%%%%
%%%             DESCRICOES TEORICAS E FERRAMENTAS BASICAS               %%%
%%%%%%%%%%%%%%%%%%%%%%%%%%%%%%%%%%%%%%%%%%%%%%%%%%%%%%%%%%%%%%%%%%%%%%%%%%%

\setlength{\parskip}{0.3cm}

\chapter{Fundamentação Teórica}
\label{ch:fundamentacao}

A cientista da computação Jeannette M. Wing, introduziu formalmente o conceito
de Pensamento Computacional (PC) em seu artigo \textit{Computational Thinking},
publicado em março de 2006, na revista \textit{Communications of the ACM}. Nesse
trabalho, \cite{wing2006} defende que o PC constitui uma habilidade cognitiva
fundamental para todos, não apenas para profissionais da área de computação. É o
processo de formular e resolver problemas.

A autora argumenta que o PC envolve abstrair, decompor problemas complexos,
reconhecer padrões e automatizar soluções, competências que, segundo ela,
deveriam integrar o currículo básico de educação ao lado da leitura, da escrita
e da aritmética.O processo de abstração consiste em filtrar os aspectos
importantes. A decomposição,em conjunto, é dividir uma tarefa grande e separar
por interesse, facilitando o gerenciamento. Isso permite a criação de soluções
de forma algorítmica. Também é descrito por \cite{wing2006} o papel do
reconhecimento de padrões para investigar e processar vastas quantidades de
dados em busca de similaridades.

Segundo \cite{wing2006} o PC é uma habilidade inerente a todos, não apenas para
cientistas da computação. É tão indispensável quanto a leitura, a escrita e a
aritmética. É pensar como os humanos resolvem problemas. 

A ciência da computação não se resume a programar, pois pensar como um cientista
da computação vai muito além de apenas ter a habilidade técnica de escrever
códigos. Na verdade, essa forma de pensar exige a capacidade de raciocinar sobre
a aplicação de conceitos computacionais para abordar e resolver problemas. É
raciocinar em múltiplos níveis de abstração.\cite{wing2006}

Nesse contexto, O PC será uma realidade quando for tão integral aos
esforços humanos que desapareça como uma filosofia explícita. Da mesma forma
que não precisamos nos lembrar conscientemente das regras gramaticais a cada
frase que dizemos, o PC deverá se integrar de forma tão
orgânica ao nosso raciocínio que deixaremos de notá-lo como uma disciplina
acadêmica isolada. É enfatizado a naturalização do PC em aspectos do
cotidiano, assim como palavras como algoritmo e pré-condição se incorporam ao
vocabulário diário de qualquer indivíduo.\cite{wing2006}

Um paralelo é a computação ubíqua, entendida como o sonho de que haveria
hardwares e softwares por todos os lados, já se tornou uma realidade, a ponto de
as pessoas não refletirem mais sobre o fato de estarem cercadas por microchips.
O pensamento computacional seguirá o mesmo caminho no futuro.\cite{wing2006}

\begin{citacao}
Não são apenas os artefatos de software e hardware que produzimos que estarão
fisicamente presentes em todos os lugares e impactará nossas vidas o tempo todo,
serão os conceitos computacionais que usamos para abordar e resolver problemas,
gerenciar nossas vidas diárias e nos comunicar e interagir com outras pessoas; e
para todos, em todos os lugares. \cite[tradução nossa]{wing2006}.
\end{citacao}

O PC, segundo \cite{wing2006}, associa-se com o raciocínio matemático e ao
pensamento de engenharia. A matemática fornece os fundamentos teóricos e formais
para a ciência da computação, permitindo formular soluções precisas para
problemas complexos. A computação exige a construção de sistemas que interagem
com o mundo real, o pensamento de engenharia torna-se indispensável para lidar
com as restrições físicas e operacionais dos dispositivos. É exatamente a
exigência de projetar ferramentas dentro dessas limitações práticas das
máquinas, somada à liberdade de construir mundos virtuais sem as mesmas amarras,
que força a transição de uma abordagem puramente matemática para o pensamento
computacional.\cite{wing2006}

A BNCC prioriza a incorporação do PC nas aulas de matemática como uma competência
essencial para a modelagem e resolução de problemas \cite{brasil2017}. Dentre
elas, destacam-se a decomposição de desafios complexos em partes menores e
gerenciáveis, o reconhecimento de padrões, além das práticas de generalização e
reúso, permitindo adaptar ou aproveitar uma lógica para resolver novos problemas
de forma genérica.

\begin{citacao}
    Decomposição é uma das principais técnicas de resolução de problemas, na
    qual um problema é dividido em subproblemas, os quais são resolvidos
    independentemente, e cujas soluções são combinadas para construir a solução
    do problema original. [...] Em conjunto, deve-se desenvolver o
    reconhecimento e a identificação de padrões, construindo conjuntos de
    objetos com base em diferentes critérios [...]. O próximo passo é comparar
    diferentes casos particulares (instâncias) de um mesmo problema,
    identificando as semelhanças e diferenças entre eles, e criar um algoritmo
    para resolver todos [...] de forma genérica. [...] Ao explicar para alguém
    como realizar uma tarefa (resolver um problema), se está criando um
    algoritmo. Esses algoritmos podem ser construídos [...] partindo do zero, na
    qual esses passos também devem ser determinados, além da sequência desses
    \cite{brasil2017}
\end{citacao}

É orientado pela BNCC o uso de tecnologias digitais como recurso pedagógico,
reconhecendo os jogos digitais como instrumentos para o desenvolvimento
de competências cognitivas. As práticas plugadas (uso direto de dispositivos e softwares) são incentivadas
como estratégias para tornar o processo de aprendizagem mais dinâmico,
significativo e engajado, aproximando os conteúdos curriculares da realidade digital
vivenciada pelos estudantes.\cite{brasil2017}

Práticas plugadas, especificamente Jogos Sérios, têm um propósito específico
relacionado ao treinamento, não apenas à diversão, como descrito por
\cite{kiryakova2014gamification}. Eles possuem todos os elementos dos jogos,
parecem jogos, mas seu objetivo é alcançar algo predeterminado.

Segundo, \cite{michael2006seriousgame} o histórico dos jogos educacionais
reflete a constante adaptação da educação a novas mídias, evoluindo de
simulações de tabuleiro na década de 1960 (como o jogo Grand Strategy) para o
grande boom do "edutainment" (educação com entretenimento) nos anos 1980 e 1990,
impulsionado pela popularização dos computadores pessoais.

Para conseguir engajar os alunos de forma eficaz, especialmente a chamada
"Geração Videogame", utiliza-se elementos interativos fundamentais do design de
jogos, associados ao conceito de gamificação, tais como o
incentivo ao aprendizado por tentativa e erro extensiva (onde o custo do
fracasso é mínimo), a entrega de feedback crítico constante, o fornecimento de
recompensas imediatas para decisões corretas. Elementos como tentativa, erro e
feedback imediato também oferecem um ambiente onde o aluno é recompensado imediatamente
por decisões corretas e, caso falhe, o erro se torna evidente no mesmo instante,
o que permite uma correção rápida.\cite{michael2006seriousgame}

O uso de elementos de gamificação com foco em melhorar o engajamento deve
englobar dimensões relacionadas à motivação, objetivos, características da
atividade e do usuário, além de aspectos emocionais e cognitivos. Segundo
\cite{silpasuwanchai2016developing}, estratégias típicas de jogos no ambiente
educacional (pontos, medalhas, desafios e quadros de líderes) têm como objetivo
elevar a motivação e, principalmente, o engajamento.

Embora o aumento do esforço e da participação tenha uma correlação positiva com
a melhoria nas notas, a percepção de diversão pelo aluno (engajamento emocional)
não se traduz automaticamente em resultados acadêmicos mais altos. A gamificação
não se deve focar apenas em entreter, mas sim em estimular um investimento
psicológico profundo que conduza à verdadeira aquisição e transferência de
habilidades para outros contextos.\cite{silpasuwanchai2016developing}

De acordo com \cite{silpasuwanchai2016developing}, os três elementos de maior
destaque, que formam a estratégia conjunta mais comum e que foram ativamente
testados no experimento prático do artigo, compõem a sigla PBL (Points, Badges,
and Leaderboards). Esses elementos podem ser aplicados como: pontuações simples,
pontos de experiência, recompensas, medalhas e distintivos concedidos por
conquistas específicas e quadros de líderes, classificando os participantes. A
adoção do PBL eleva de forma significativa o engajamento comportamental e
cognitivo dos estudantes, ampliando a atenção visual, o esforço e a persistência
voluntária em atividades acadêmicas, favorecendo a evolução na resolução de
problemas ao longo do tempo. 

O emprego do PBL é adaptável, pois atua como um incentivo fundamental para reter
alunos em tarefas consideradas monótonas ou difíceis  e serve como um forte
combustível motivacional para estudantes com perfil focado no sucesso
\textit{(high-achievers)}, que se dedicam intensamente ao desafio impulsionados
pelo desejo de superar os demais e dominar os
placares.\cite{silpasuwanchai2016developing}
